\chapter{予測と機械学習}\label{ch:ml}

この章では、「何かを学習させて予測する」ときの基本と、その落とし穴を説明する。この研究で学習させたのは、リッジ回帰、FinBERT の追加学習、ChronoGPT の LoRA、ロジスティック回帰である（\chapref{ch:io}）。どれも、この章の決まりに従って学習と評価を分けた。

\section{教師あり学習}

入力 $x$（文章のベクトル、数字の特徴など）から出力 $y$（反応、上がったかどうか）を予測する関数 $f$ を、正解のわかっているデータ $\{(x_i,y_i)\}$ から作ることを\term{教師あり学習}と呼ぶ。$y$ が連続の値なら\term{回帰}、カテゴリなら\term{分類}である。

学習に使ったデータ（\term{学習データ}）でよく当たっても意味はない。大事なのは、学習に使っていない新しいデータ（\term{テストデータ}）での当たり方、つまり\term{汎化}の性能である。

\section{過学習と正則化}

\subsection{過学習}

モデルが複雑すぎる（調整できるパラメータが多すぎる）と、学習データの雑音まで覚えてしまい、新しいデータでは外れる。これを\term{過学習}と呼ぶ。決算の反応は雑音がとても大きい（点数で説明できるのは数 \% だけ）ので、過学習しやすい。

\begin{exbox}
ChronoGPT の埋め込みは 1 発表あたり 1,536 個の数である。学習の最初の時期（2015 年の 1 年分の発表だけで学習する時期）には、データの数が係数の数と同じ程度しかない。普通の最小二乗法では、学習データをほぼ完全に当てる係数が見つかってしまい、新しいデータではほとんど当たらない。
\end{exbox}

\subsection{リッジ回帰}

\term{リッジ回帰}は、係数の大きさに罰則を付けて過学習を抑える。
\begin{equation}
\hat\beta_{\text{ridge}} = \arg\min_\beta \;\|y - X\beta\|^2 + \lambda\|\beta\|^2 = (X\T X + \lambda I)^{-1}X\T y
\end{equation}
$\lambda$ が大きいほど、係数は 0 に向けて縮む。$X\T X$ が特異でも、$\lambda I$ を足すので逆行列が存在する。共分散の縮小推定（\chapref{ch:cov}）と同じ発想である。

$\lambda$ は\term{交差検証}で選ぶ。学習データをいくつかに分け、一部で学習して残りで当て、当たりが最もよい $\lambda$ を使う。リッジでは、1 つずつ抜いた交差検証が行列の計算 1 回でできる（\code{RidgeCV}）。実験 5 では、どの四半期でも $\lambda$ は $1.8\times10^4$〜$5.6\times10^4$ が選ばれ、候補の端（$10^5$）には当たらなかった。

\subsection{特徴の標準化}

罰則は係数の大きさに掛かるので、特徴の単位が違うと罰則の効き方が変わる。学習の前に、各特徴を学習データの平均と標準偏差で標準化する。このとき、\textbf{テストデータの平均や標準偏差を使ってはならない}（未来の情報が混ざる）。

\section{時間のあるデータの学習：ローリング}

株価のデータには時間の順序がある。ランダムに学習とテストに分けると、未来のデータで学習して過去を当てることになる。そこで\term{ローリング}（ウォーク・フォワード）で学習する。

\begin{figure}[htbp]
\centering
\begin{tikzpicture}[x=9mm,y=6mm]
\foreach \r/\a/\b in {0/0/4,1/0/5,2/0/6,3/0/7}{
  \fill[c1!25] (\a,-\r) rectangle (\b,-\r+0.7);
  \fill[c2!35] (\b,-\r) rectangle (\b+1,-\r+0.7);
  \node[font=\scriptsize,anchor=east] at (-0.2,-\r+0.35) {四半期 \pgfmathparse{int(\b+1)}\pgfmathresult};
}
\node[font=\scriptsize] at (2,1.2) {学習（反応が確定した過去の発表）};
\node[font=\scriptsize,color=c2!80!black] at (7.2,1.2) {予測};
\draw[->,color=ink2] (0,-3.9) -- (9,-3.9) node[right,font=\scriptsize] {時間};
\end{tikzpicture}
\caption{ローリングの学習。四半期ごとに、それより前に反応が確定した発表だけで学び直し、その四半期の発表を予測する。}
\label{fig:ml-rolling}
\end{figure}

\begin{defbox}[この研究の決まり]
四半期 $q$ の発表を予測するモデルは、\textbf{四半期 $q$ が始まる前に反応の窓（day 1）が閉じた発表だけ}で学習する。反応の窓が四半期をまたぐ発表は、学習に入れない。
\end{defbox}

\section{ロジスティック回帰と AUC}

\subsection{ロジスティック回帰}

$y$ が 0 か 1（反応が正か負か）のとき、
\begin{equation}
\Pr(y=1\mid x) = \frac{1}{1+\exp(-x\T\beta)}
\end{equation}
とおき、対数尤度を最大にする $\beta$ を求めるのが\term{ロジスティック回帰}である。係数は「$x_j$ が 1 増えると、オッズ $p/(1-p)$ が $e^{\beta_j}$ 倍になる」と読める。

\subsection{AUC}

\term{AUC}（ROC 曲線の下の面積）は、「ランダムに選んだ $y=1$ の例と $y=0$ の例で、モデルが $y=1$ の例により高い確率を付ける確率」である。0.5 は当てずっぽう、1 は完全である。株価の上げ下げの予測では、0.55〜0.6 でも意味がある。

\begin{notebox}[外れ値に弱い]
実験 11 で事前登録したロジスティック回帰（数字の特徴をそのまま入れる）は、AUC 0.593 と、アナリスト予想との差だけの単純な順位（0.612）にも負けた。驚きの値には極端な外れ値があり、ロジスティック回帰の係数がそれに引っぱられたためである。特徴を順位に直した版（探索）は 0.611 になった。外れ値の多い金融のデータでは、順位にする・丸めるといった前処理が効く。
\end{notebox}

\section{データの漏れ}

学習や評価に、予測の時点では知りえない情報が入り込むことを\term{漏れ}（リーク）と呼ぶ。漏れがあると、テストの成績が実際より良く見える。

\begin{longtable}{P{4.0cm}P{10.4cm}}
\toprule
\textbf{漏れの種類} & \textbf{例と、この研究での防ぎ方} \\
\midrule
\endhead
\bottomrule
\endfoot
時間の漏れ & 未来の発表で学習する。→ ローリングで、反応の窓が閉じた過去だけで学習 \\
前処理の漏れ & 全期間の平均で標準化する。→ 学習データの統計量だけを使う \\
目的変数の漏れ & 未来の実現リターンで $\Omega$ を決める。→ 禁止（実験の約束） \\
モデルの知識の漏れ & LLM が学習のときに、その後の株価やニュースを読んでいる。→ 知識の締め切りより後の期間だけを使う（gpt-4o-mini は 2023-11 以降）。または版の決まった ChronoGPT を使う（\chapref{ch:pit}） \\
データの時点の漏れ & 後から訂正された財務の数値を、当時の値として使う。→ 最初に提出された値を使う（総合報告書 1 の実験 C） \\
選び方の漏れ & 結果を見てから、良い設定を選ぶ。→ 事前登録と、開発・最終の分離 \\
\end{longtable}

\section{事前登録}

\term{事前登録}は、実験の計画（問い、データ、変数、検定、判定の規則）を、\textbf{結果を見る前に}文書にして固定することである。医学の臨床試験で始まり、心理学や経済学に広がった。

\begin{keybox}[この研究の事前登録のやり方]
\begin{enumerate}
\item 計画の文書（\code{preregistration.md}）と評価のコード（\code{run.py}）を書く。株価の反応は一度も計算しない。
\item 2 つを\textbf{同じコミット}で Git に固定する。
\item コードを 1 回だけ実行する。実行に使ったコミットは結果のファイルに記録される。
\item 文書に書いていない分析は、すべて「探索」として分けて報告する。
\item 計画からずれたら（コードの誤りなど）、そのことと、直した版の結果の両方を報告する（実験 9 の日付、実験 11 の修正 1 の規則）。
\end{enumerate}
\end{keybox}

事前登録の目的は、「分析の分かれ道」（どの変数を入れるか、どの期間にするか、外れ値をどう扱うか）を、結果に合わせて選べなくすることである。分かれ道が 10 個あって 2 通りずつなら、1,024 通りの分析ができる。その中から有意なものを選べば、効果がなくても有意な結果を報告できてしまう。

\subsection{開発と最終の分離}

修正を\textbf{探す}実験（実験 11）では、事前登録だけでは足りない。何を試すかの候補が多いからである。そこで、
\begin{enumerate}
\item データを時間で\term{開発}の期間と\term{最終}の期間に分ける（実験 11 の主な切り方 B では、開発は 2016〜2019 年、最終は 2020 年以降）。
\item 開発の期間だけで、事前に決めた機械的な規則（差の $t$ が 1 以上、前半・後半とも悪くならない、似ていれば単純なほう）で修正を選ぶ。人は選ばない。
\item 選ばれた修正を、最終の期間で 1 回だけ評価する。
\end{enumerate}
とした。最終の期間は、選ぶ過程で一度も使わないので、その結果は選び方の偏りを含まない。

\section{人工データで水増しできるか}

金融のデータは少ない（決算は 1 社年 4 回）。深層学習には足りないので、「生成したデータで水増しすればよいのでは」という考えがある。López de Prado は、過学習を\textbf{検出する}ために人工のデータ（既知の性質を持つ価格の経路）を使うことを勧めている。しかし、\textbf{学習のデータを増やす}目的には、次の理由で使えない。

\begin{itemize}
\item 人工データは、元のデータ（またはその仮定のモデル）から作られる。元のデータにない「文章と反応の関係」の情報は、人工データにも入っていない。
\item 元のデータに過学習させた生成モデルで作ったデータは、元のデータの雑音をそのまま写す。それで学習すると、雑音をさらに強く覚える。
\end{itemize}

情報を増やすには、新しい本物のデータ（期間を延ばす、別の種類の文章を足す）か、別の課題で学習した知識（事前学習したモデル）を借りるしかない。後者が、ChronoGPT や FinBERT を使う理由である（\chapref{ch:llm}）。

\section{深層学習の基本}

\term{ニューラルネットワーク}は、線形の変換と非線形の関数を何層も重ねたモデルである。パラメータ（重み）は、予測の誤差（\term{損失}）を小さくする向きに少しずつ動かして学習する（\term{勾配降下法}）。

\begin{description}[leftmargin=1.5em,style=nextline]
\item[バッチ] 1 回の更新に使うデータのまとまり。
\item[歩（ステップ）] 1 回の更新。
\item[エポック] 学習データ全体を 1 回なめること。実験 4 は 2 エポック（1,748 歩）、実験 10 は 1 エポック（873 歩）。
\item[学習率] 1 歩で重みを動かす大きさ。
\item[損失の推移] 学習が進むと損失が下がる。実験 10 の損失は、目的変数を標準化しているので、何も学ばなければ約 1.0 である。実際に約 1.0 からほとんど下がらなかった。
\end{description}

\begin{projbox}
\begin{itemize}
\item \textbf{リッジ}：TF-IDF（実験 2、$\lambda=10$ に固定）と ChronoGPT の埋め込み（実験 5・9・11、交差検証）。四半期ごとのローリング。
\item \textbf{ロジスティック回帰}：実験 11 の最後の比較。数字だけと、数字＋点数の AUC の差。
\item \textbf{深層学習}：FinBERT の全体の追加学習（実験 4・11）、ChronoGPT の LoRA（実験 10）。
\item \textbf{事前登録}：実験 2〜11、X02〜X04。探索は X01・X05 と、各実験の「探索」の節。
\item \textbf{開発と最終の分離}：実験 11（切り方 A・B・C）。
\end{itemize}
\end{projbox}
